\section{Results}
\label{sec:results}

\subsection{The headline ratio}
For the sourced lunar-regolith seed, \eqref{eq:copy} gives a copy time of about 582~days,
so the derived dwell of \eqref{eq:dwell} is $\tau \approx 1.59$~years. At fleet scale this
same 582~days is the doubling clock: neighbours are days apart, so build time is the only
clock that matters. Run as the settlement-front dwell for a powered cruise across a field of
1200 Sun-like stars, it sits against a fill time $T_{100} \approx 1.9\times10^{6}$~years,
reproducing to the same order of magnitude the multi-million-year powered timescales of the
model we build on~\cite{nicholson-forgan-2013}. The ratio of \eqref{eq:frac} is then
\begin{equation}
f \;\approx\; 8.5\times10^{-7},
\end{equation}
under one part in a million. The constraint that is the entire clock at fleet scale is, at
galactic scale, a rounding error.

\subsection{The result survives not knowing the copy time}
The copy time is a lunar-factory proxy, so the case cannot rest on its exact value; it rests
on the margin. Fig.~\ref{fig:margin} sweeps the copy time across five orders of magnitude and
records the resulting dwell fraction of the powered fill. The fraction rises only in step with
the copy time, and stays below one percent of the fill until the copy time is about
$1.5\times10^{4}$ times its nominal value, that is, a copy time near 25\,000 years rather than
582 days. A hypothetical probe whose true build cadence were four orders of magnitude slower
than the lunar proxy would still leave the verdict unchanged. The conclusion is therefore
insensitive to the one number we cannot yet source, which is the strongest statement we can
make given the gap. The documented uncertainty in the collector inputs to \eqref{eq:copy}, of
order a factor of two, is a small sub-interval of this margin.

\begin{figure}[t]
\centering
\includegraphics[width=\columnwidth]{fig_margin.pdf}
\caption{The manufacturing dwell as a fraction of the powered galactic fill time, swept over
five orders of magnitude in copy time (log-log; 1200-star field, event-resolved). The nominal
lunar-seed copy time (582~days, dotted line) gives a fraction under $10^{-6}$. The fraction
reaches the one-percent bar (dashed) only at a copy time about $1.5\times10^{4}$ times nominal
(square), so an error of four orders of magnitude in the build cadence does not change the
verdict.}
\label{fig:margin}
\end{figure}

\subsection{The same holds across travel policies}
The ratio is smallest for the slowest front and largest for the fastest, because a faster
fill leaves the fixed dwell a larger share of a smaller total. Table~\ref{tab:policy} sweeps
the three named policies on a common field and reports, for each, the break-even multiplier:
the factor by which the copy time would have to grow before the dwell reached one percent of
that policy's fill. Even for the fastest front, nearest-target slingshots, the copy time would
have to be about sixty times larger, a probe taking a century rather than a year and a half to
build one copy, before manufacturing became a one-percent effect. The fastest speeds
contemplated in the literature, beamed directed-energy cruise at 0.1 to 0.2~$c$~\cite{lubin-2016},
lie beyond the slingshot regime we simulate and would push the crossover nearer still; the
break-even multipliers bound how much nearer, and none of the simulated policies comes within
a factor of sixty of it.

\begin{table}[t]
\caption{Dwell fraction and break-even copy-time multiplier by travel policy, on a common
400-star field (event-resolved). The break-even multiplier is the factor by which the copy
time must grow before the dwell reaches one percent of that policy's fill time.}
\label{tab:policy}
\centering
\begin{tabular}{lccc}
\toprule
Policy & Fill time (yr) & Dwell fraction & Break-even ($\times$ nominal) \\
\midrule
Powered            & $1.4\times10^{6}$ & $1.2\times10^{-6}$ & $\sim$10\,000 \\
Slingshot max-boost& $4.8\times10^{4}$ & $3.3\times10^{-5}$ & $\sim$590 \\
Slingshot nearest  & $8.3\times10^{3}$ & $1.9\times10^{-4}$ & $\sim$60 \\
\bottomrule
\end{tabular}
\end{table}

\subsection{Measuring the tax directly}
The ratio $f$ is an analytic separation; we also measure the dwell's cost on the front
directly, as the fractional increase in fill time when the derived dwell is switched on
against the old zero-dwell baseline on the same field. For a multi-million-year powered fill
this difference is far too small to resolve by simulation, which is itself the finding, so we
measure it for the fast slingshot policies, where it is largest, over a 24-seed ensemble
(Fig.~\ref{fig:tax}). For nearest-target slingshots the tax is tight and clearly resolved at a
median of 0.32~percent of exploration time (interquartile range 0.26 to 0.36~percent). For
max-boost slingshots the per-seed values scatter from about $-6$ to $+10$~percent around a
median near one percent: at this field size the dwell's effect is within the seed-to-seed
noise of the fill time, so we report it as positive-but-unresolved rather than as a point
estimate. Reporting a single seed here, as an earlier draft of the underlying finding did,
would have hidden that spread.

\begin{figure}[t]
\centering
\includegraphics[width=\columnwidth]{fig_dwell_tax.pdf}
\caption{The measured dwell tax (fractional increase in fill time from switching the derived
dwell on) per seed for the two slingshot policies, 24 seeds each, event-resolved on a 400-star
field. Bars are the ensemble median, shaded bands the interquartile range. Nearest-target
slingshots give a tight, resolved 0.32~percent; max-boost scatters across zero and is within
seed noise.}
\label{fig:tax}
\end{figure}

Two checks confirm the measurement is not an artifact. Halving the integration timestep leaves
the nearest-slingshot tax stable and converging on the event-resolved value, whereas a timestep
coarse against the dwell under-resolves it, the same discretization hazard identified for the
underlying front model; we therefore report event-resolved values throughout. Varying the field
size over 200, 400, and 800 stars moves the median tax only between 0.31 and 0.34~percent, so
the small-field measurement is representative over the range tested.
